\documentclass[11pt,a4paper]{article}
\usepackage[T1]{fontenc}
\usepackage[utf8]{inputenc}
\usepackage{lmodern}
\usepackage[a4paper,margin=27mm]{geometry}
\usepackage{amsmath,amssymb,amsthm,mathtools}
\usepackage{mathrsfs}
\usepackage{microtype}
\usepackage{enumitem}
\usepackage{xcolor}
\usepackage{tikz}
\usetikzlibrary{arrows.meta,positioning,calc,patterns}
\setlist{itemsep=3pt,topsep=5pt}
\setlength{\emergencystretch}{2em}
\numberwithin{equation}{section}

\newtheorem{theorem}{Theorem}[section]
\newtheorem{proposition}[theorem]{Proposition}
\newtheorem{lemma}[theorem]{Lemma}
\newtheorem{corollary}[theorem]{Corollary}
\newtheorem{maintheorem}{Theorem}
\renewcommand{\themaintheorem}{\Alph{maintheorem}}
\theoremstyle{definition}
\newtheorem{definition}[theorem]{Definition}
\theoremstyle{remark}
\newtheorem{remark}[theorem]{Remark}

\definecolor{fblue}{RGB}{42,95,155}
\definecolor{forange}{RGB}{180,98,35}
\definecolor{fgreen}{RGB}{35,130,120}
\definecolor{fpurple}{RGB}{125,47,145}
\definecolor{fred}{RGB}{170,40,40}

\newcommand{\R}{\mathbb R}
\newcommand{\C}{\mathbb C}
\newcommand{\Z}{\mathbb Z}
\newcommand{\N}{\mathbb N}
\newcommand{\T}{\mathbb T}
\newcommand{\Tf}{\mathbb T_\eta}
\newcommand{\ip}[2]{\langle #1,#2\rangle}
\newcommand{\one}{\mathbf 1}
\newcommand{\dd}{\mathrm d}
\newcommand{\ii}{\mathrm i}
\newcommand{\eps}{\varepsilon}
\newcommand{\cK}{\mathcal K}
\newcommand{\cF}{\mathcal F}
\newcommand{\cG}{\mathcal G}
\newcommand{\cV}{\mathcal V}
\newcommand{\cW}{\mathcal W}
\newcommand{\cX}{\mathcal X}
\newcommand{\cD}{\mathcal D}
\newcommand{\cN}{\mathcal N}
\newcommand{\bQ}{\mathbf Q}
\newcommand{\bG}{\mathbf G}
\newcommand{\bH}{\mathbf H}
\newcommand{\bC}{\mathbf C}
\newcommand{\bD}{\mathbf D}
\newcommand{\bL}{\boldsymbol{\Delta}}
\newcommand{\Lap}{\Delta}
\newcommand{\dl}{d_1}
\newcommand{\Hol}[2]{[#1]_{\alpha,#2}}
\newcommand{\Lam}{\Lambda}
\newcommand{\tn}[1]{|\!|\!|#1|\!|\!|}
\newcommand{\fA}{\mathfrak A}
\DeclareMathOperator{\supp}{supp}
\DeclareMathOperator{\ran}{ran}
\DeclareMathOperator{\Rea}{Re}



% --- Details volume additions ---
\usepackage{booktabs,tabularx,longtable}
\usepackage{xr-hyper}
\usepackage[colorlinks=true,linkcolor=blue!50!black,citecolor=blue!50!black,urlcolor=blue!50!black]{hyperref}
\externaldocument[N-]{../main}[../YM4-small-field-flow.pdf]   % links go to the shipped note
\newcommand{\Nref}[1]{\ref{N-#1}}
\newcommand{\Neqref}[1]{\eqref{N-#1}}
\newtheorem{step}{Step}[section]
\newtheorem{detail}[theorem]{Detail}
% Steps are numbered within each Detail.
\counterwithin*{step}{theorem}\renewcommand{\thestep}{\arabic{step}}
\makeatletter
% Named destination N:<note label>, emitted once (the first Detail for a label wins).
\newcommand{\Ntarget}[1]{\@for\RL@t:=#1\do{\ifcsname RL@tgt@\RL@t\endcsname\else
  \expandafter\gdef\csname RL@tgt@\RL@t\endcsname{}%
  \edef\RL@x{\noexpand\hypertarget{N:\RL@t}{}}\RL@x\fi}}
\makeatother
% A step that opens a detail: start the detail's paragraph first, so that the step's link anchor is emitted once.
\makeatletter\AddToHook{env/step/before}{\if@inlabel\leavevmode\fi}\makeatother
\newcommand{\status}[1]{\par\noindent\textbf{Status:} #1\par}
% Names of table (ii) of the note (as in ../rl-preamble.tex).
\newcommand{\Cav}{C^{\rm av}}
\newcommand{\Av}{\operatorname{Av}}
\newcommand{\cJ}{\mathcal J}
\newcommand{\Eloop}{E_{\rm loop}}
